\section{Introduction}
Foundation models for biology such as BioCLIP~\citep{stevens2024bioclip} classify images of organisms zero-shot at
any rank of the taxonomy. They are integrated into citizen-science tools, camera-trap pipelines and the detection of
invasive species, where a confident wrong name has a cost. Standard benchmarks measure accuracy on curated photos;
they say little about the images such a tool meets in practice: unusual backgrounds, night flash, larvae instead of
adults, partial views, or deliberately misleading pictures.

We study this question with generative models. Our application produces \emph{semantic} adversarial examples: images
that a person (and an independent model) recognises as species $y$, but on which BioCLIP outputs another species.
Unlike pixel-level attacks~\citep{szegedy2014intriguing}, these images are plausible photographs; they are in the
spirit of unrestricted~\citep{song2018constructing} and natural~\citep{hendrycks2021natural} adversarial examples.

Our definition makes the problem well posed: an image only counts as an attack if an evaluation oracle that is
independent of the training signal confirms the species (Section~\ref{sec:method}). Without this condition, a generator
can trivially ``fool'' the classifier by drawing badly; an RL-trained generator learns exactly that (we show it).

\textbf{Contributions.}
(i) A pipeline that attacks the real, open-set BioCLIP with text-to-image and image-to-image generation, on four
datasets.
(ii) A post-training method: DDPO~\citep{black2024ddpo} with group-relative advantages~\citep{shao2024deepseekmath}
and a two-oracle reward, which generalises to species unseen in training.
(iii) An inference-time improvement, \emph{early rejection}, that prunes best-of-$N$ candidates from a cheap preview
of the predicted clean image, and a study of sampling acceleration.
(iv) An analysis of \emph{where} BioCLIP fails: per context family, per taxonomic rank, in its embedding space and in
its saliency maps. Figure~\ref{fig:pipeline} summarises the approach.

\begin{figure}[t]
\centering
\resizebox{\columnwidth}{!}{% Summary figure of the report. \usepackage{tikz} \usetikzlibrary{positioning,arrows.meta,fit}
\begin{tikzpicture}[
  box/.style={draw, rounded corners, align=center, font=\small, minimum height=8mm, inner sep=4pt},
  gen/.style={box, fill=blue!10}, clf/.style={box, fill=orange!15}, opt/.style={box, fill=green!12},
  arr/.style={-{Latex[length=2mm]}, thick}]
  \node[box] (sp) {species $y$\\(80, 4 groups)};
  \node[box, right=6mm of sp] (pr) {prompt builder\\modifiers / targeted / LLM};
  \node[gen, right=6mm of pr] (sd) {Stable Diffusion 1.5\\+ LoRA $\phi$};
  \node[clf, above right=4mm and 10mm of sd] (t) {BioCLIP (target)\\$p_T(y)$};
  \node[clf, right=10mm of sd] (o) {BioCLIP-2 (reward oracle)\\$p_O(y)$};
  \node[clf, below right=4mm and 10mm of sd, fill=gray!15] (e) {VLM judge + humans\\(evaluation only)};
  \node[opt, below=14mm of pr, minimum width=60mm] (r) {$r=(1-p_T(y))\,p_O(y)$ \quad best-of-$N$ / early rejection / RAFT / DDPO (GRPO)};
  \draw[arr] (sp) -- (pr);
  \draw[arr] (pr) -- (sd);
  \draw[arr] (sd) -- (t);
  \draw[arr] (sd) -- (o);
  \draw[arr, dashed] (sd) -- (e);
  \draw[arr] (t.east) -| ++(4mm,-30mm) |- (r.east);
  \draw[arr] (o.east) -| ++(2mm,-14mm) |- (r.east);
  \draw[arr] (r.north) -- node[right, font=\scriptsize]{update $\phi$ / select} (sd.south);
\end{tikzpicture}
}
\caption{BioFool. A prompt (or a real photo) is turned into an image by Stable Diffusion with a LoRA. BioCLIP (the
target) and BioCLIP-2 (the reward oracle) score it; the reward drives search and RL. Reported fidelity comes from a
separate VLM judge and from humans.}
\label{fig:pipeline}
\end{figure}
